\section{Conclusão}

O desenvolvimento deste projeto permitiu a consolidação prática de conceitos fundamentais de sistemas digitais sequenciais, por meio de tentativa e erro, incluindo a construção de circuitos combinacionais, a síntese de funções booleanas com uso de Mapas de Karnaugh e a compreensão de Máquinas de Estados Finitas.

A abordagem estritamente baseada em componentes básicos exigiu a decomposição explícita da lógica interna do sistema e compreensão de alguns conceitos físicos que permeiam esta área do conhecimento. Outrossim, evidenciou o papel de elementos como multiplexadores na seleção entre modos operacionais de incremento e decremento, bem como a diferenciação de propagações síncronas e assíncronas.

A modelagem do sistema como uma arquitetura modular baseada em FSMs encadeadas permitiu uma separação clara entre os níveis de controle e armazenamento de estado, reforçando a importância da organização hierárquica em sistemas sequenciais complexos.

Como etapa complementar, a prototipagem no ambiente Minecraft Java Edition mostrou-se útil para a validação conceitual do sistema, permitindo observar na prática limitações relacionadas à propagação de sinais, sincronização e atrasos lógicos, frequentemente abstraídos em ferramentas de simulação tradicionais.

A utilização de flip-flops do tipo D como elemento básico de armazenamento confirmou a viabilidade da construção de sistemas temporais determinísticos a partir de unidades discretas de memória sincronizadas por clock global.

Por fim, o trabalho demonstra que a construção de sistemas digitais complexos depende não apenas da correta definição de blocos lógicos individuais, mas principalmente da coerência na integração entre camadas de abstração, onde a estrutura de estados e a dinâmica de transição desempenham papel central no comportamento global do sistema.